\clearpage
\begin{proofbox}
\par\medskip
\noindent\begin{minipage}{\linewidth}
\paragraph{证明一：最直接的想法——两头等高，中间总有一处走平}
\begin{center}% F081: precise mathematical teaching diagram.
\begin{tikzpicture}[x=1cm,y=1cm,>=stealth,font=\small]
\draw[rounded corners=3pt,gray!40] (0,0) rectangle (14.2,9.0);
\begin{scope}[yshift=1cm]
\node at (7.1,7.4) {内部最高点：两侧割线逼向同一条水平切线};
\begin{scope}[shift={(4.2,2.5)},x=2.3cm,y=2cm]
\draw[->,gray!65] (-1.35,0)--(1.35,0) node[right] {$x$};
\draw[->,gray!65] (0,-0.3)--(0,1.45) node[above] {$y$};
\draw[blue!75!black,line width=1.15pt,,domain=-1:1,samples=160] plot (\x,{1-(\x)*(\x)});
\fill[blue!75!black] (-1,0) circle (1.8pt);\node[below left] at (-1,0) {$A$};
\fill[blue!75!black] (1,0) circle (1.8pt);\node[below right] at (1,0) {$B$};
\fill[blue!75!black] (0,1) circle (1.8pt);\node[above left] at (0,1) {$C$};
\draw[orange!85!black,line width=1.2pt] (-0.35,1)--(0.35,1);
\draw[green!50!black,line width=1pt] (-.7,.51)--(0,1);\draw[purple,line width=1pt] (0,1)--(.7,.51);
\fill[green!50!black] (-0.7,0.51) circle (1.8pt);\node[above left] at (-0.7,0.51) {$h<0$};
\fill[purple] (0.7,0.51) circle (1.8pt);\node[above right] at (0.7,0.51) {$h>0$};
\draw[dashed,gray!60] (0,0)--(0,1);\node[below right] at (0,0) {$c$};
\end{scope}
\node[green!50!black] at (10.4,5.3) {左侧差商 $\ge0$};
\node[purple] at (10.4,4.3) {右侧差商 $\le0$};
\node[] at (10.4,3.2) {可导：两侧极限相同};
\node[] at (10.4,2.1) {所以 $f^{\prime}(c)=0$};
\end{scope}
\node[text width=13cm,align=center] at (7.1,.9) {闭区间连续保证能取到最高值；两端等高且有更高点，最高点就在内部。};
\end{tikzpicture}
\end{center}
\end{minipage}\par\medskip

用闭区间连续函数的最大、最小值及导数定义，给出\textbf{一般证明}。

\textbf{第1步：先想最简单的情况。}
如果从 $a$ 到 $b$ 的函数值一直不变，那么中间任一点的导数都是 $0$，结论成立。

\textbf{第2步：如果中间不是一直保持同一个高度。}
令两端的共同高度为 $L$。由闭区间连续性，函数能取得最大值和最小值。若内部有点高于 $L$，最大值就大于两端的值，故在内部取得；若有点低于 $L$，同理最小值在内部取得。不恒为 $L$ 时，两种情况至少有一种发生。记相应位置为 $c$，则 $a<c<b$。

\textbf{第3步：为什么这个位置的导数等于 $0$？}
先看 $c$ 是最高位置。取很小的 $h$，只要 $c+h$ 仍在区间里，就有
\[
f(c+h)-f(c)\le0.
\]
因此：
\[
\begin{array}{ll}
h>0:&\displaystyle\frac{f(c+h)-f(c)}{h}\le0,\\[6pt]
h<0:&\displaystyle\frac{f(c+h)-f(c)}{h}\ge0.
\end{array}
\]
\textbf{因为在 $c$ 可导，两侧差商趋于同一个数。这个数既不小于 $0$ 又不大于 $0$，只能等于 $0$。}

于是 $f'(c)=0$。如果 $c$ 是最低位置，不等号的方向反过来，同样得到 $f'(c)=0$。

\textbf{把刚才的证明翻译成几何语言：}两端的连线斜率是
\[
\frac{f(b)-f(a)}{b-a}=0;
\]
证明已经找到一个内部位置 $c$，它的切线斜率 $f'(c)$ 也为 $0$。\textbf{所以在这两点之间，至少有一条切线与两端连线平行。}这是罗尔定理要表达的图形关系。

\textbf{结论：}恒定或不恒定的情形都能找到内部的 $c$，满足 $f'(c)=0$。
\end{proofbox}
\clearpage
\begin{proofbox}
\par\medskip
\noindent\begin{minipage}{\linewidth}
\paragraph{证明二：用初中二次函数亲手验算——看清水平切线在哪里}
\begin{center}% F081: precise mathematical teaching diagram.
\begin{tikzpicture}[x=1cm,y=1cm,>=stealth,font=\small]
\draw[rounded corners=3pt,gray!40] (0,0) rectangle (14.2,9.0);
\begin{scope}[yshift=1cm]
\node at (7.1,7.4) {二次函数验算：零点的中点就是最低点};
\begin{scope}[shift={(3.5,6)},x=1.55cm,y=0.8cm]
\draw[->,gray!65] (-1.4,0)--(3.45,0) node[right] {$x$};
\draw[->,gray!65] (0,-4.8)--(0,0.7) node[above] {$y$};
\draw[blue!75!black,line width=1.15pt,,domain=-1:3,samples=160] plot (\x,{(\x+1)*(\x-3)});
\fill[blue!75!black] (-1,0) circle (1.8pt);\node[above left] at (-1,0) {$a=-1$};
\fill[blue!75!black] (3,0) circle (1.8pt);\node[above right] at (3,0) {$b=3$};
\fill[blue!75!black] (1,-4) circle (1.8pt);\node[below] at (1,-4) {$c=1$};
\draw[orange!85!black,line width=1.2pt] (0.35,-4)--(1.65,-4);
\end{scope}
\node[] at (10.8,3.5) {$f^{\prime}(c)=0$};
\end{scope}
\node[text width=13cm,align=center] at (7.1,.9) {图取 $a=-1,b=3$；一般式 $(x-a)(x-b)$ 在 $c=(a+b)/2$ 处导数为零。};
\end{tikzpicture}
\end{center}
\end{minipage}\par\medskip

这是一种\textbf{特别容易看懂的具体例子}，用来帮助理解证明一；它并不是对所有函数的另一份一般证明。

取任意两个数 $a<b$，构造
\[
f(x)=(x-a)(x-b).
\]
很容易看出：
\[
f(a)=0,\qquad f(b)=0.
\]
把式子展开：
\[
f(x)=x^2-(a+b)x+ab.
\]
求导得到：
\[
f'(x)=2x-(a+b).
\]
令导数等于 $0$，得到
\[
\boxed{\ c=\frac{a+b}{2}\ }.
\]
因为 $a<b$，所以
\[
a<\frac{a+b}{2}<b.
\]
它恰好位于两端的\textbf{中间}。从图像上看，这就是抛物线的最低位置，切线刚好水平。

\textbf{别误会：}一般函数中的 $c$ \textbf{不一定是两个端点的中点}。本例的对称性保证了中点处导数为零，不能把它当作一般结论。
\end{proofbox}
\clearpage
\begin{proofbox}
\par\medskip
\noindent\begin{minipage}{\linewidth}
\paragraph{证明三：不用求出原方程的解，也能知道它最多有几个}
\begin{center}% F081: precise mathematical teaching diagram.
\begin{tikzpicture}[x=1cm,y=1cm,>=stealth,font=\small]
\draw[rounded corners=3pt,gray!40] (0,0) rectangle (14.2,9.2);
\begin{scope}[yshift=1cm]
\node at (7.1,7.6) {四个不同零点，夹出三个不同的导数零点};
\begin{scope}[shift={(7.1,4.2)},x=1.65cm,y=0.16cm]
\draw[->,gray!65] (-3.6,0)--(3.6,0) node[right] {$x$};
\draw[->,gray!65] (0,-19)--(0,14) node[above] {$y$};
\draw[blue!75!black,line width=1.15pt,,domain=-3.12:3.12,samples=160] plot (\x,{(\x)*(\x)*(\x)*(\x)-10*(\x)*(\x)+9});
\fill[blue!75!black] (-3,0) circle (1.8pt);\node[above right] at (-3,0) {$x_1$};
\fill[blue!75!black] (-1,0) circle (1.8pt);\node[above left] at (-1,0) {$x_2$};
\fill[blue!75!black] (1,0) circle (1.8pt);\node[above right] at (1,0) {$x_3$};
\fill[blue!75!black] (3,0) circle (1.8pt);\node[above left] at (3,0) {$x_4$};
\draw[orange!85!black,line width=1.2pt] (-2.5560679774997896,-16.000000000000004)--(-1.9160679774997897,-16.000000000000004);
\fill[orange!85!black] (-2.23606797749979,-16.000000000000004) circle (1.8pt);\node[below] at (-2.23606797749979,-16.000000000000004) {$c_1$};
\draw[orange!85!black,line width=1.2pt] (-0.32,9)--(0.32,9);
\fill[orange!85!black] (0,9) circle (1.8pt);\node[above left] at (0,9) {$c_2$};
\draw[orange!85!black,line width=1.2pt] (1.9160679774997897,-16.000000000000004)--(2.5560679774997896,-16.000000000000004);
\fill[orange!85!black] (2.23606797749979,-16.000000000000004) circle (1.8pt);\node[below] at (2.23606797749979,-16.000000000000004) {$c_3$};
\end{scope}
\node[] at (7.1,.65) {$x_1<c_1<x_2<c_2<x_3<c_3<x_4$};
\end{scope}
\node[text width=13cm,align=center] at (7.1,.9) {图例 $f(x)=(x^2-1)(x^2-9)$；每个开区间各给一个，位置不会重复。};
\end{tikzpicture}
\end{center}
\end{minipage}\par\medskip

现在证明\textbf{零点个数与导数零点个数的关系}。

\textbf{第1步：}假设 $f(x)$ 有四个不同零点：
\[
x_1<x_2<x_3<x_4.
\]
\textbf{第2步：}分别在三对相邻零点上使用罗尔定理：
\[
\begin{aligned}
(x_1,x_2)&\text{ 内有 }c_1,\quad f'(c_1)=0,\\
(x_2,x_3)&\text{ 内有 }c_2,\quad f'(c_2)=0,\\
(x_3,x_4)&\text{ 内有 }c_3,\quad f'(c_3)=0.
\end{aligned}
\]
它们分别位于三个\textbf{互不重叠}的小区间，因此 $c_1,c_2,c_3$ 是三个不同的数。

\textbf{第3步：}所以只要原函数有四个不同零点，导数就至少要有三个不同零点。如果我们已经算出导数\textbf{最多只有两个零点}，那么原函数有四个零点就不可能发生。

同理，若导数最多有 $k$ 个零点，原函数最多只能有 $k+1$ 个不同零点。

\textbf{再进一步：数与任意水平直线的交点。}若想知道图像与 $y=h$ 有几个交点，只需把方程写成
\[
 f(x)=h\quad\Longleftrightarrow\quad g(x)=f(x)-h=0.
\]
因为减去一个固定数字不会改变导数，$g'(x)=f'(x)$，所以刚才的推理原样成立：\textbf{水平切线位置最多有 $k$ 个，图像与任何一条水平直线的不同交点就最多有 $k+1$ 个}。



\textbf{还可以再往下数一次：}
\begin{center}% F081: precise mathematical teaching diagram.
\begin{tikzpicture}[x=1cm,y=1cm,>=stealth,font=\small]
\draw[rounded corners=3pt,gray!40] (0,0) rectangle (14.2,8.6);
\begin{scope}[yshift=1cm]
\node at (7.1,7.0) {再用一次：三个零点逐层夹出两个、一个};
\draw[->,gray!70] (3,5.9)--(12.3,5.9);
\node[] at (1.55,5.9) {$f=0$};
\fill[blue!75!black] (3.5999999999999996,5.9) circle (2pt);\node[] at (3.5999999999999996,6.300000000000001) {$x_1$};
\fill[blue!75!black] (7.1,5.9) circle (2pt);\node[] at (7.1,6.300000000000001) {$x_2$};
\fill[blue!75!black] (10.6,5.9) circle (2pt);\node[] at (10.6,6.300000000000001) {$x_3$};
\draw[->,gray!70] (3,3.9)--(12.3,3.9);
\node[] at (1.55,3.9) {$f^{\prime}=0$};
\fill[blue!75!black] (5.07927405783631,3.9) circle (2pt);\node[] at (5.07927405783631,4.3) {$u$};
\fill[blue!75!black] (9.12072594216369,3.9) circle (2pt);\node[] at (9.12072594216369,4.3) {$v$};
\draw[->,gray!70] (3,1.9)--(12.3,1.9);
\node[] at (1.55,1.9) {$f^{\prime\prime}=0$};
\fill[blue!75!black] (7.1,1.9) circle (2pt);\node[] at (7.1,2.3) {$w$};
\draw[dashed,orange!75!black] (3.5999999999999996,5.75)--(5.07927405783631,4.55);
\draw[dashed,orange!75!black] (7.1,5.75)--(5.07927405783631,4.55);
\draw[dashed,orange!75!black] (7.1,5.75)--(9.12072594216369,4.55);
\draw[dashed,orange!75!black] (10.6,5.75)--(9.12072594216369,4.55);
\draw[dashed,orange!75!black] (5.07927405783631,3.75)--(7.1,2.55);
\draw[dashed,orange!75!black] (9.12072594216369,3.75)--(7.1,2.55);
\end{scope}
\node[text width=13cm,align=center] at (7.1,.9) {图取 $f=x^3-x$：$u=-1/\sqrt3,\ v=1/\sqrt3,\ w=0$。};
\end{tikzpicture}
\end{center}
如果 $f(x)$ 有三个不同零点，那么 $f'(x)$ 至少有两个不同零点；接着对 $f'(x)$ 再用一次罗尔定理，就知道 $f''(x)$ 至少有一个零点。

这里设 $f$ 在同一个区间 $I$ 上连续，在其内部二阶可导（即内部每一点都能求出 $f^{\prime\prime}$）。若原函数有三个零点，第一次夹出的 $u,v$ 都在 $I$ 内部；$f^{\prime}$ 在 $[u,v]$ 上连续，在 $(u,v)$ 内可导，可以再次使用罗尔定理。因此，
\[
\boxed{\ f''(x)\text{ 在区间内没有零点}\ \Longrightarrow\ f(x)\text{ 在该区间内最多有两个不同零点}\ }.
\]
\textbf{这个方法的用处：}有些原方程很难解，导数或再求一次导数却很简单；这时“反过来数”比硬解方程快得多。
\end{proofbox}


% -------- 新增验证：九种基本组合都有具体的函数 --------
\clearpage
\begin{proofbox}
\paragraph{补充验证：九种走向真的都能出现吗？}

\textbf{只用一个容易算的式子，就能造出全部九种。}令
\[
 f(x)=
 \begin{cases}
   ax^3,&x<0,\\
   bx^3,&x\ge 0,
 \end{cases}
 \qquad a,b\text{ 各从 }-1,0,1\text{ 中任选一个}.
\]
\textbf{第一步：检查中间的切线是否水平。}无论 $a,b$ 怎样取，$f(0)=0$，而
\[
\frac{f(h)-f(0)}{h}=
\begin{cases}ah^2,&h<0,\\bh^2,&h>0.\end{cases}
\]
当 $h$ 越来越接近 $0$ 时，两侧的值都越来越接近 $0$，所以\textbf{$f'(0)=0$}。

\textbf{第二步：分别看左右怎样走。}在 $x\ne0$ 时，
\[
 f'(x)=\begin{cases}3ax^2,&x<0,\\3bx^2,&x>0.\end{cases}
\]
因为 $x^2>0$，左边导数的正负由 $a$ 决定，右边由 $b$ 决定：
$a,b$ 取正数就上升，取 $0$ 就平着走，取负数就下降。
这样选 $a,b$，\textbf{九种组合全部可以实现}。

\textbf{要点：}九宫格是\textbf{左右各自保持一种导数符号}时的九种基本走法。
如果某一侧的导数还会反复变号，这一侧就不能简单标成一个``升''或``降''，需要继续分段研究。
\end{proofbox}
